%----------------------------------------------------------------------------------------------------------------------------
\chapter*{Prefazione}
\addcontentsline{toc}{chapter}{Prefazione}
%----------------------------------------------------------------------------------------------------------------------------

Questi tre volumi nascono dall'esperienza didattica maturata negli anni di insegnamento della Scienza delle Costruzioni
presso la Facoltà di Ingegneria dell'\textit{Uni\-versità degli	Studi di Roma ``La Sapienza''}. La trasformazione da appunti
di lezione a testo scritto non è una semplice trascrizione: ha richiesto di rendere esplicita l'architettura logica che
nei corsi si trasmette anche attraverso la voce, il gesto e il dialogo con gli studenti, e che sulla pagina deve reggersi
da sola.

\section*{I tre volumi e i corsi}

I tre volumi sono pensati come percorso consecutivo --- il
secondo presuppone gli argomenti del primo, il terzo quelli
dei due precedenti --- ma ciascuno è autocontenuto: chi
padroneggi i prerequisiti del volume scelto può leggerlo
come unità singola.

Il rapporto fra i tre volumi e l'insegnamento è articolato. Il
primo volume è pensato per gli studenti della laurea
triennale che seguono il corso di Scienza delle Costruzioni;
il secondo volume copre la parte successiva del medesimo
corso ed estende la trattazione alla visione nodale dei
sistemi di travi e alla trave su suolo elastico; il terzo
volume mescola argomenti del corso triennale --- analisi
della capacità portante, carico critico, verifiche di
sicurezza --- con argomenti che il primo anno della laurea
magistrale dedica alle strutture bidimensionali, ovvero le
travi curve e le funi. Questi ultimi capitoli (trave su
suolo elastico, travi curve, funi) costituiscono propriamente
la prima fase di un corso di meccanica delle strutture
bidimensionali, e vi sono inclusi perché forniscono al
lettore le chiavi per interpretare il funzionamento delle
strutture bidimensionali piane e curve: sono il passaggio
concettuale fra il modello monodimensionale della trave e i
modelli di superficie. La parte di Vol.~II e Vol.~III che
eccede i corsi triennali non è dunque un addendum: è ciò che
permette ai tre volumi di essere consultati anche oltre la
durata di un singolo corso.

Nel loro insieme i tre volumi coprono in modo
sostanzialmente completo la meccanica della trave in regime
lineare e in alcune sue estensioni non lineari classiche
(analisi limite, stabilità).

Il dettaglio dei contenuti di ciascun volume è esposto nella
\PARref{sec:i_tre_volumi} del capitolo introduttivo.



\section*{A chi è rivolto questo volume}

Questo primo volume è scritto per lo studente che affronta lo
studio sistematico della meccanica della trave e dei sistemi
di travi nel modello rigido. I prerequisiti sono ridotti a un
nucleo essenziale: gli strumenti matematici di base appresi
nei corsi del primo anno --- calcolo differenziale e
integrale, geometria analitica nel piano e nello spazio,
nozioni elementari di algebra lineare --- e una familiarità
con la meccanica del corpo rigido nei suoi termini fisici:
corpi, forze, vincoli, equilibrio. Tutti gli strumenti
specifici --- algebra delle matrici e sistemi lineari,
applicazioni lineari e dualità, vettori geometrici nelle loro
varietà --- sono ripresi nella Parte~I (\textit{Preliminari})
con il livello di dettaglio e l'orientamento richiesti dalle
applicazioni strutturali. Il volume è quindi accessibile a
chi affronta la disciplina per la prima volta, sia studente
della laurea triennale in Ingegneria o in Architettura, sia
lettore in autonomia che voglia rifondare in modo organico la
propria comprensione del modello rigido.

L'obiettivo non è fornire un ricettario di formule da
applicare, ma costruire una comprensione profonda dei
meccanismi fisici e della struttura logica della teoria. Chi
avrà interiorizzato questa struttura --- la dualità fra
cinematica e statica, la classificazione dei sistemi rigidi
secondo i quattro casi dell'algebra lineare, il principio dei
lavori virtuali come strumento operativo unitario --- sarà in
grado di orientarsi in contesti nuovi, di riconoscere la
stessa architettura in modelli diversi, di usare gli
strumenti computazionali con consapevolezza critica anziché
meccanicamente.


\tG{\section*{Il debito intellettuale}

Ogni testo di meccanica delle strutture si inserisce in una tradizione. Questa non fa eccezione.

Giulio Ceradini ha insegnato a generazioni di ingegneri romani che la meccanica delle strutture non è una collezione
di tecniche ma una teoria unitaria, fondata su principi precisi e capace di produrre risultati generali. Il suo modo
di impostare i problemi --- partire sempre dai fondamenti, non dare nulla per scontato, cercare la struttura logica
prima del risultato numerico --- è il modello a cui questo volume aspira. Questo modo di fare meccanica si è trasmesso
attraverso i suoi allievi --- Vincenzo Ciampi, Antonio Di Carlo, Carlo Gavarini, Angelo Luongo, Giuseppe Rega,
Fabrizio Vestroni --- e ha plasmato una generazione di ricercatori e docenti romani. Di quella tradizione, e dei
suoi scritti, queste pagine si riconoscono debitrici.}


\section*{La struttura dei tre volumi}

I tre volumi costituiscono un percorso unitario che ha per
protagonista la trave --- il modello strutturale più
semplice e al tempo stesso il più ricco di conseguenze. Il
primo volume si articola in tre parti. La prima
(\emph{Preliminari}) raccoglie gli strumenti matematici e
concettuali su cui poggia l'intera trattazione: algebra
delle matrici e sistemi lineari, applicazioni lineari e
dualità, vettori geometrici --- liberi, applicati, polari e
assiali. La seconda parte sviluppa la \emph{cinematica e la
	statica del corpo rigido e dei sistemi di corpi rigidi} in
piena generalità, attraverso gli strumenti fondamentali
della disciplina: classificazione cinematica e statica,
dualità cinematica-statica, principio dei lavori virtuali.
La terza parte specializza la trattazione al caso della
\emph{trave e dei sistemi di travi rigide} --- travature a
collegamento eterogeneo e travature a collegamento omogeneo,
queste ultime nelle loro due famiglie classiche dei sistemi
reticolari e dei telai piani. Il volume si chiude mettendo
in evidenza le indeterminazioni statiche e le impossibilità
cinematiche che il modello rigido non è in grado di
risolvere, motivando l'introduzione della deformabilità.

Il secondo volume applica questa grammatica alla trave
rettilinea deformabile. Si apre con un capitolo di richiami
della Meccanica dei Solidi --- cinematica della deformazione,
equazioni di equilibrio di Cauchy, legame costitutivo
elastico lineare --- che costituisce il ponte verso la
derivazione delle equazioni della trave: i diagrammi di
Tonti, presentati sistematicamente in un'appendice e
utilizzati come strumento di lettura ricorrente, rendono
visibile l'architettura logica della meccanica delle strutture
e mostrano come le equazioni della trave si ottengano dal
continuo tridimensionale per riduzione dimensionale. Si
sviluppano la teoria della trave secondo i modelli di
Eulero-Bernoulli e Timoshenko, le caratteristiche della
sollecitazione e i metodi di analisi delle strutture
isostatiche e iperstatiche, fino ai sistemi di travi
deformabili nelle loro due letture --- per assemblaggi e
nodale --- e alla trave su suolo elastico.

Il terzo volume mette la grammatica alla prova, estendendola
in direzioni che ne rivelano la portata: la trave curva e la
fune, le non linearità costitutive e geometriche, le
verifiche di sicurezza strutturale.



Al termine di questo percorso lo studente possiede gli strumenti per affrontare qualsiasi teoria strutturale basata su ipotesi cinematiche. La stessa grammatica si ritrova immutata nei modelli bidimensionali --- lastre, piastre, membrane e gusci --- che costituiscono la naturale prosecuzione nei corsi avanzati. Non ci sarà nulla da reimparare da zero: solo la stessa architettura in nuove vesti.

\section*{Un invito alla lettura}

La meccanica delle strutture è un edificio logico di rara coerenza, costruito in oltre due secoli di lavoro scientifico da menti straordinarie. Studiarlo con attenzione --- non limitandosi a seguire i calcoli ma cercando di capire perché ogni passaggio è necessario, quale principio fisico lo giustifica, quale struttura logica lo organizza --- è un'esperienza intellettuale di prim'ordine.

Questi volumi cercano di rendere accessibile questa esperienza. Le formule ci sono, i calcoli ci sono, gli esempi ci sono: ma l'obiettivo finale non è saper calcolare le travi. È acquisire quella \textit{sensibilità strutturale} --- la capacità di intuire il comportamento di una struttura prima ancora di calcolarla, di riconoscere le soluzioni efficaci da quelle problematiche, di anticipare i punti critici --- che distingue l'ingegnere strutturista dal calcolatore.

Questa sensibilità si costruisce nel tempo e attraverso la pratica. Ma ha le sue radici nella comprensione profonda dei fondamenti. È lì che questi volumi cercano di portare il lettore.


\bigskip

\noindent Roma, \the\year

\noindent \textit{Gli Autori}


%%----------------------------------------------------------------------------------------------------------------------------
%\chapter*{Prefazione}
%\addcontentsline{toc}{chapter}{Prefazione}
%%----------------------------------------------------------------------------------------------------------------------------
%
%Questi tre volumi nascono dall'esperienza didattica maturata negli anni di insegnamento della Scienza delle Costruzioni
%presso la Facoltà di Ingegneria dell'\textit{Uni\-versità degli Studi di Roma ``La Sapienza''}. La trasformazione da appunti
%di lezione a testo scritto non è una semplice trascrizione: ha richiesto di rendere esplicita l'architettura logica che
%nei corsi si trasmette anche attraverso la voce, il gesto e il dialogo con gli studenti, e che sulla pagina deve reggersi
%da sola.
%
%
%\section*{Struttura dell'opera}
%
%I tre volumi costituiscono un percorso unitario che ha per protagonista la trave --- il modello strutturale più
%semplice e al tempo stesso il più ricco di conseguenze. Sono pensati come percorso consecutivo --- il secondo
%presuppone gli argomenti del primo, il terzo quelli dei due precedenti --- ma ciascuno è autocontenuto: chi padroneggi
%i prerequisiti del volume scelto può leggerlo come unità singola.
%
%Il primo volume si apre con una parte preliminare (\emph{Preliminari}) che raccoglie gli strumenti matematici e
%concettuali necessari: algebra delle matrici e sistemi lineari, applicazioni lineari e dualità, vettori geometrici
%nelle loro varietà di liberi, applicati, polari e assiali. Sviluppa quindi la cinematica e la statica del corpo rigido
%e dei sistemi di corpi rigidi, e ne specializza la trattazione al caso della trave e dei sistemi di travi rigide ---
%a collegamento eterogeneo e a collegamento omogeneo, questi ultimi nelle loro due famiglie classiche dei sistemi
%reticolari e dei telai piani. Si chiude mettendo in evidenza le indeterminazioni statiche e le impossibilità
%cinematiche che il modello rigido non è in grado di risolvere, motivando l'introduzione della deformabilità. Il
%secondo volume applica la stessa grammatica alla trave rettilinea deformabile, dai modelli di Eulero-Bernoulli e
%Timoshenko ai sistemi di travi deformabili nelle loro due letture --- per assemblaggi e nodale --- fino alla trave
%su suolo elastico. Il terzo volume mette la grammatica alla prova estendendola in direzioni che ne rivelano la
%portata: la trave curva e la fune, le non linearità costitutive e geometriche, le verifiche di sicurezza strutturale.
%
%Il rapporto fra l'opera e l'insegnamento è articolato. Il primo volume è pensato per gli studenti della laurea
%triennale che seguono il corso di Scienza delle Costruzioni; il secondo copre la parte successiva del medesimo corso.
%Il terzo mescola argomenti del corso triennale --- analisi della capacità portante, carico critico, verifiche di
%sicurezza --- con argomenti che il primo anno della laurea magistrale dedica alle strutture bidimensionali, ovvero
%la trave curva e la fune. Questi ultimi capitoli (trave su suolo elastico, trave curva, fune) costituiscono
%propriamente la prima fase di un corso di meccanica delle strutture bidimensionali, e vi sono inclusi perché
%forniscono al lettore le chiavi per interpretare il funzionamento delle strutture bidimensionali piane e curve:
%sono il passaggio concettuale fra il modello monodimensionale della trave e i modelli di superficie. La parte di
%Vol.~II e Vol.~III che eccede i corsi triennali non è dunque un addendum: è la parte che conferisce all'opera il
%suo carattere di riferimento ragionato, oltre la durata di un singolo corso.
%
%Nel loro insieme i tre volumi coprono in modo sostanzialmente completo la meccanica della trave in regime lineare
%e in alcune sue estensioni non lineari classiche (analisi limite, stabilità). La stessa grammatica si ritrova
%immutata nei modelli bidimensionali --- lastre, piastre, membrane, gusci --- che costituiscono la naturale
%prosecuzione nei corsi avanzati: non c'è nulla da reimparare da zero, solo la stessa architettura in nuove vesti.
%
%Il dettaglio dei contenuti di ciascun volume è esposto nel \PARref{sec:i_tre_volumi} del capitolo introduttivo.
%
%
%\section*{A chi è rivolto questo volume}
%
%Questo primo volume è scritto per lo studente che affronta lo studio sistematico della meccanica della trave e dei
%sistemi di travi nel modello rigido. I prerequisiti sono ridotti a un nucleo essenziale: gli strumenti matematici di
%base appresi nei corsi del primo anno --- calcolo differenziale e integrale, geometria analitica nel piano e nello
%spazio, nozioni elementari di algebra lineare --- e una familiarità con la meccanica del corpo rigido nei suoi
%termini fisici: corpi, forze, vincoli, equilibrio. Tutti gli strumenti specifici sono ripresi nella Parte~I
%(\emph{Preliminari}) con il livello di dettaglio e l'orientamento richiesti dalle applicazioni strutturali. Il volume
%è quindi accessibile a chi affronta la disciplina per la prima volta, sia studente della laurea triennale in
%Ingegneria o in Architettura, sia lettore in autonomia che voglia rifondare in modo organico la propria comprensione
%del modello rigido.
%
%L'architettura del volume si regge su tre cardini che ne attraversano l'intera trattazione: la dualità fra cinematica
%e statica, la classificazione dei sistemi rigidi secondo i quattro casi dell'algebra lineare, il principio dei lavori
%virtuali come strumento operativo unitario. Chi avrà interiorizzato questi cardini saprà ritrovarli in contesti
%nuovi e usare gli strumenti computazionali con consapevolezza critica anziché meccanicamente.
%
%
%\section*{Un invito alla lettura}
%
%La meccanica delle strutture è un edificio logico di rara coerenza, costruito in oltre due secoli di lavoro
%scientifico da menti straordinarie. Studiarlo con attenzione --- non limitandosi a seguire i calcoli ma cercando di
%capire perché ogni passaggio è necessario, quale principio fisico lo giustifica, quale struttura logica lo organizza
%--- è un'esperienza intellettuale di prim'ordine.
%
%Questi volumi cercano di rendere accessibile quell'esperienza. Le formule ci sono, i calcoli ci sono, gli esempi ci
%sono: ma l'obiettivo finale non è saper calcolare le travi. È acquisire quella \textit{sensibilità strutturale} ---
%la capacità di intuire il comportamento di una struttura prima ancora di calcolarla, di riconoscere le soluzioni
%efficaci da quelle problematiche, di anticipare i punti critici --- che distingue l'ingegnere strutturista dal
%calcolatore.
%
%Questa sensibilità si costruisce nel tempo e attraverso la pratica. Ma ha le sue radici nella comprensione profonda
%dei fondamenti. È lì che questi volumi cercano di portare il lettore.
%
%
%\section*{Il debito intellettuale}
%
%Ogni testo di meccanica delle strutture si inserisce in una tradizione. Questo non fa eccezione.
%
%Giulio Ceradini ha insegnato a generazioni di ingegneri romani che la meccanica delle strutture non è una collezione
%di tecniche ma una teoria unitaria, fondata su principi precisi e capace di produrre risultati generali. Il suo modo
%di impostare i problemi --- partire sempre dai fondamenti, non dare nulla per scontato, cercare la struttura logica
%prima del risultato numerico --- è il modello a cui questo volume aspira. Questo modo di fare meccanica si è
%trasmesso attraverso i suoi allievi --- Vincenzo Ciampi, Antonio Di Carlo, Carlo Gavarini, Angelo Luongo, Giuseppe
%Rega, Fabrizio Vestroni --- e ha plasmato una generazione di ricercatori e docenti romani; è attraverso questa catena
%di trasmissione, oltre che attraverso i suoi scritti, che la sua impronta è presente in queste pagine.
%
%
%\bigskip
%
%\noindent Roma, \the\year
%
%\noindent \textit{Gli Autori}